\section{A1 for \texorpdfstring{$\Spin(2n{+}1)$}{Spin(2n+1)}: explicit one-plaquette coefficient bounds}
\label{app:A1_Spin_odd}

This appendix proves \Cref{prop:A1_Spin_odd}.
The structure parallels Appendix~\ref{app:A1_Sp}, with the only genuine differences being:
(i) the maximal-torus trace has an extra constant eigenvalue $1$,
(ii) the Weyl shift is half-integral (type $B_n$),
(iii) the resulting moment kernel has a $+1$ Hankel shift.

\subsection{Maximal torus and Wilson weight}
\label{subsec:A1_Spin_odd_setup}

Fix $n\ge 2$ and set $G=\Spin(2n+1)$.
Since the Wilson weight factors through the projection to $\SO(2n+1)$, only the sector $\pi(-\mathbf 1)=\mathbf 1$ contributes
(\Cref{rem:center_vanish}).
We therefore work on $\SO(2n+1)$.

Every conjugacy class meets the maximal torus consisting of block-diagonal rotations
with eigenvalues
\begin{equation}
\label{eq:SO_odd_eigs}
\{1,e^{\pm i\theta_1},\dots,e^{\pm i\theta_n}\},\qquad \theta_j\in[0,\pi].
\end{equation}
In the defining $(2n+1)$-dimensional representation $\rho$,
\begin{equation}
\label{eq:SO_odd_trace}
\Re\Tr_{\rho}(U)=1+2\sum_{j=1}^n \cos\theta_j.
\end{equation}
The constant term $1$ only contributes a global factor to the one-plaquette weight.
Thus, up to normalization,
\Cref{eq:one_plaq_measure_general} corresponds on the torus to
\begin{equation}
\label{eq:SO_odd_symbol}
W_t(\theta):=\exp\Big(t\sum_{j=1}^n \cos\theta_j\Big),
\qquad t:=\frac{2\alpha\beta}{2n+1}.
\end{equation}

\subsection{Type \texorpdfstring{$B_n$}{B\_n} Weyl character formula and Andr\'eief reduction}
\label{subsec:A1_Spin_odd_determinant}

Irreducible representations of $\SO(2n+1)$ are indexed by dominant integral highest weights,
which we realize as partitions $\lambda=(\lambda_1\ge\cdots\ge\lambda_n\ge 0)$.
Define integer shifts
\begin{equation}
\label{eq:SO_odd_bi_int}
 b_i:=n-i\qquad(1\le i\le n),
\end{equation}
so that the type-$B_n$ Weyl vector has components
\begin{equation}
\label{eq:SO_odd_rho}
\rho_i=b_i+\tfrac12=n-i+\tfrac12.
\end{equation}
Define the shifted integer indices
\begin{equation}
\label{eq:SO_odd_Ai}
A_i(\lambda):=\lambda_i+b_i=\lambda_i+n-i\in\N_0.
\end{equation}

\begin{lemma}[Weyl character formula for $\SO(2n{+}1)$ in half-integer sine form]
\label{lem:SOodd_character_sine}
Let $U\in T$ have eigenangles $\theta=(\theta_1,\dots,\theta_n)$.
Then
\begin{equation}
\label{eq:SOodd_character_formula}
\chi_{\lambda}(\theta)
=\frac{\det\Big(\sin\big((A_i(\lambda)+\tfrac12)\theta_j\big)\Big)_{1\le i,j\le n}}
        {\det\Big(\sin\big((b_i+\tfrac12)\theta_j\big)\Big)_{1\le i,j\le n}}.
\end{equation}
\end{lemma}

\begin{lemma}[Weyl integration formula in determinant form for $\SO(2n{+}1)$]
\label{lem:SOodd_weyl_integration}
There exists a constant $c_{B}(n)>0$ such that for every continuous class function $F$ on $\SO(2n+1)$,
\begin{equation}
\label{eq:SOodd_weyl_integration}
\int_{\SO(2n+1)}F(U)\,d\mu_{\mathrm{Haar}}(U)
= c_{B}(n)\int_{[0,\pi]^n}F(\theta)\,\det\Big(\sin\big((b_i+\tfrac12)\theta_j\big)\Big)^2\,d\theta.
\end{equation}
\end{lemma}

\begin{proof}
This is the Weyl integration formula for type $B_n$ written using the determinant representation of the Weyl denominator.
\end{proof}

\begin{proposition}[Determinantal formula for the one-plaquette coefficients on $\Spin(2n{+}1)$]
\label{prop:SOodd_det_formula}
Let $\lambda$ be a partition of length at most $n$.
Define the half-integer sine moment kernel
\begin{equation}
\label{eq:SOodd_kernel_def}
M_t^{B}(p,q)
:=\frac{2}{\pi}\int_0^{\pi} e^{t\cos\theta}\,\sin\big((p+\tfrac12)\theta\big)\,\sin\big((q+\tfrac12)\theta\big)\,d\theta,
\qquad p,q\in\N_0.
\end{equation}
Let
\begin{equation}
\label{eq:SOodd_Dlambda_def}
D_{\lambda}^{B}(t)
:=\det\Big(M_t^{B}\big(A_i(\lambda),\,b_j\big)\Big)_{1\le i,j\le n}.
\end{equation}
Then with $t=2\alpha\beta/(2n+1)$ one has
\begin{equation}
\label{eq:SOodd_coeff_det_ratio}
 a_{\lambda}(\beta,\alpha)
=\E_{\nu_{\beta,\alpha}}\Big[\frac{\chi_{\lambda}(U)}{\dim\lambda}\Big]
=\frac{1}{\dim\lambda}\,\frac{D_{\lambda}^{B}(t)}{D_{\varnothing}^{B}(t)}.
\end{equation}
Moreover, for every $p,q\in\N_0$,
\begin{equation}
\label{eq:SOodd_kernel_bessel}
M_t^{B}(p,q)=I_{|p-q|}(t)-I_{p+q+1}(t).
\end{equation}
\end{proposition}

\begin{proof}
The proof is identical to the $\Sp(2n)$ case (Appendix~\ref{app:A1_Sp}) after replacing the type-$C$ determinant formulas
by Lemmas~\ref{lem:SOodd_character_sine} and \ref{lem:SOodd_weyl_integration}.
The cancellation of one Weyl-denominator factor and Andr\'eief's identity produce \Cref{eq:SOodd_Dlambda_def}.

To compute the kernel, use
\(\sin((p+\tfrac12)\theta)\sin((q+\tfrac12)\theta)=\tfrac12(\cos((p-q)\theta)-\cos((p+q+1)\theta))\)
so that
\begin{align*}
M_t^{B}(p,q)
&=\frac{1}{\pi}\int_0^{\pi} e^{t\cos\theta}\,\big(\cos((p-q)\theta)-\cos((p+q+1)\theta)\big)\,d\theta
\\
&=I_{|p-q|}(t)-I_{p+q+1}(t).
\end{align*}
\end{proof}

\subsection{One-box inequality and telescoping}
\label{subsec:A1_Spin_odd_one_box}

\begin{lemma}[One-box multiplicative bound for type $B_n$]
\label{lem:SOodd_one_box}
Fix $t>0$.
Let $\mu$ be a partition of length at most $n$ and let $i$ be an index such that $\mu+e_i$ is still a partition.
Write $A_i:=A_i(\mu)=\mu_i+b_i$.
Then
\begin{equation}
\label{eq:SOodd_one_box_bound}
\frac{a_{\mu+e_i}}{a_{\mu}}\ \le\ \frac{I_{A_i+1}(t)}{I_{A_i}(t)}.
\end{equation}
\end{lemma}

\begin{proof}
\textbf{Step 1: TN row shift bound for determinants.}
By Lemma~\ref{lem:moment_matrix_TN}(ii) the matrix $(M_t^{B}(p,q))_{p,q\ge 0}$ is TN.
Apply Lemma~\ref{lem:TN_row_shift} with the increasing column tuple $C=(0,1,\dots,n-1)$
(i.e. $b_n<b_{n-1}<\cdots<b_1$).
As in Appendix~\ref{app:A1_Sp} this yields
\begin{equation}
\label{eq:SOodd_det_shift}
\frac{D^{B}_{\mu+e_i}(t)}{D^{B}_{\mu}(t)}
\ \le\ \frac{\Delta(A(\mu+e_i))}{\Delta(A(\mu))}\,\frac{M_t^{B}(A_i+1,0)}{M_t^{B}(A_i,0)}.
\end{equation}

\textbf{Step 2: cancel the Vandermonde factor using the type-$B$ Weyl dimension ratio.}

\paragraph{Denominator-minor convention (TN ratios).}
At several points below we form Vandermonde-normalized minor ratios (or equivalent determinant ratios) and invoke the TN one-row shift tool from Appendix~\ref{app:TP_tools}.
Whenever a would-be denominator minor can vanish, we apply the zero-case clause from Lemma~J.5:
if the pivot entry (hence the denominator minor) vanishes, then the corresponding numerator determinant also vanishes and the desired inequality is trivial; otherwise we work in the strictly-positive regime and proceed with the ratio bound.



Write $a_i:=A_i+\tfrac12$ and $\rho_i:=b_i+\tfrac12$.
The type-$B_n$ Weyl dimension formula can be written as
\begin{equation}
\label{eq:SOodd_dim_formula}
\dim\lambda
=\Bigg(\prod_{i=1}^n \frac{a_i(\lambda)}{\rho_i}\Bigg)
\Bigg(\prod_{1\le i<j\le n}\frac{a_i(\lambda)^2-a_j(\lambda)^2}{\rho_i^2-\rho_j^2}\Bigg).
\end{equation}
Thus, under $\mu\mapsto\mu+e_i$ (so $a_i\mapsto a_i+1$ and $a_j$ fixed for $j\ne i$),
\begin{equation}
\label{eq:SOodd_dim_ratio}
\frac{\dim(\mu+e_i)}{\dim(\mu)}
=\frac{a_i+1}{a_i}\prod_{j\ne i}\frac{(a_i+1)^2-a_j^2}{a_i^2-a_j^2}
=\frac{a_i+1}{a_i}\prod_{j\ne i}\frac{a_i+1-a_j}{a_i-a_j}\cdot\frac{a_i+1+a_j}{a_i+a_j}.
\end{equation}
Since $a_i-a_j=A_i-A_j$, the product over $(a_i+1-a_j)/(a_i-a_j)$ is exactly the Vandermonde ratio
$\Delta(A(\mu+e_i))/\Delta(A(\mu))$.
Therefore
\begin{equation}
\label{eq:SOodd_dim_cancel}
\frac{\dim(\mu)}{\dim(\mu+e_i)}\,\frac{\Delta(A(\mu+e_i))}{\Delta(A(\mu))}
=\frac{a_i}{a_i+1}\prod_{j\ne i}\frac{a_i+a_j}{a_i+1+a_j}\ \le\ \frac{a_i}{a_i+1}\ \le\ 1.
\end{equation}

\textbf{Step 3: bound the kernel entry ratio by an adjacent Bessel ratio.}
Using \Cref{eq:SOodd_kernel_bessel} at column $0$,
\begin{equation}
\label{eq:SOodd_col0_entries}
M_t^{B}(p,0)=I_p(t)-I_{p+1}(t)\qquad(p\ge 0).
\end{equation}
Hence
\begin{equation}
\label{eq:SOodd_Mratio}
\frac{M_t^{B}(A_i+1,0)}{M_t^{B}(A_i,0)}
=\frac{I_{A_i+1}-I_{A_i+2}}{I_{A_i}-I_{A_i+1}}.
\end{equation}
By Lemma~\ref{lem:bessel_ratio_monotone}, the ratio $I_{k+1}/I_k$ is nondecreasing in $k$.
Thus $I_{A_i}I_{A_i+2}\ge I_{A_i+1}^2$, which is equivalent to
$I_{A_i}(I_{A_i+1}-I_{A_i+2})\le I_{A_i+1}(I_{A_i}-I_{A_i+1})$.
Dividing by the positive denominator $I_{A_i}(I_{A_i}-I_{A_i+1})$ yields
\begin{equation}
\label{eq:SOodd_Mratio_bound}
\frac{M_t^{B}(A_i+1,0)}{M_t^{B}(A_i,0)}\ \le\ \frac{I_{A_i+1}(t)}{I_{A_i}(t)}.
\end{equation}

\textbf{Step 4: combine.}
Using \Cref{eq:SOodd_coeff_det_ratio} we have
$\frac{a_{\mu+e_i}}{a_{\mu}}=\frac{\dim\mu}{\dim(\mu+e_i)}\cdot\frac{D^{B}_{\mu+e_i}}{D^{B}_{\mu}}$.
Combine \Cref{eq:SOodd_det_shift}, \Cref{eq:SOodd_dim_cancel}, and \Cref{eq:SOodd_Mratio_bound} to obtain
\Cref{eq:SOodd_one_box_bound}.
\end{proof}

\begin{corollary}[Product bound over boxes]
\label{cor:SOodd_telescoping}
Let $\lambda$ be a nonempty partition of length at most $n$.
Then
\begin{equation}
\label{eq:SOodd_product_over_boxes}
0\le a_{\lambda}(\beta,\alpha)
\ \le\ \prod_{(i,j)\in\lambda}\frac{I_{b_i+j}(t)}{I_{b_i+j-1}(t)},
\qquad t=\frac{2\alpha\beta}{2n+1},\quad b_i=n-i.
\end{equation}
\end{corollary}

\begin{proof}
Iterate \Cref{eq:SOodd_one_box_bound} along any one-box growth chain from $\varnothing$ to $\lambda$.
At the moment the $j$th box in row $i$ is added, $A_i=b_i+(j-1)$.
\end{proof}

\subsection{Crossover bound in Casimir form}
\label{subsec:A1_Spin_odd_crossover}

\begin{lemma}[Casimir in partition coordinates for type $B_n$]
\label{lem:SOodd_casimir}
Let $\lambda$ be a partition of length at most $n$.
With the root-length convention of \Cref{def:metric_casimir},
\begin{equation}
\label{eq:SOodd_casimir_formula}
C_2(\lambda)=\sum_{i=1}^n \lambda_i\big(\lambda_i+2b_i+1\big),\qquad b_i=n-i.
\end{equation}
\end{lemma}

\begin{proof}
For type $B_n$, $\rho_i=b_i+\tfrac12$ and $C_2(\lambda)=\langle\lambda,\lambda+2\rho\rangle$.
Expanding yields \Cref{eq:SOodd_casimir_formula}.
\end{proof}

\begin{proposition}[A1 crossover bound for $\Spin(2n{+}1)$]
\label{prop:SOodd_A1_final}
There exists an explicit constant $c_{B}(n)>0$ such that for every nontrivial $\lambda$ and every $\alpha\in[1/2,1]$, $\beta\ge 1$,
\begin{equation}
\label{eq:SOodd_A1_final}
 a_{\lambda}(\beta,\alpha)
\ \le\
\exp\Big(-c_{B}(n)\,\frac{C_2(\lambda)}{\beta+\sqrt{C_2(\lambda)}}\Big).
\end{equation}
\end{proposition}

\begin{proof}
Fix $t=2\alpha\beta/(2n+1)$.
Take logarithms of \Cref{eq:SOodd_product_over_boxes} and apply Lemma~\ref{lem:amos_upper_bound}:
\begin{equation}
\label{eq:SOodd_log_asinh}
\log a_{\lambda}\ \le\ -\sum_{(i,j)\in\lambda}\operatorname{arsinh}\Big(\frac{b_i+j-\tfrac12}{t}\Big).
\end{equation}
As in Appendix~\ref{app:A1_Sp},
$\operatorname{arsinh}(x)\ge x/(1+x)$ and $x/(1+x)\ge x/(t+K_{\max})$ give
\begin{equation}
\label{eq:SOodd_sum_lower}
\sum_{(i,j)\in\lambda}\operatorname{arsinh}\Big(\frac{b_i+j-\tfrac12}{t}\Big)
\ge \frac{\sum_{(i,j)\in\lambda}(b_i+j-\tfrac12)}{t+K_{\max}},
\end{equation}
with $K_{\max}=\max_{(i,j)\in\lambda}(b_i+j-\tfrac12)=b_1+\lambda_1-\tfrac12\le (n-1)+\lambda_1$.

Compute the numerator:
\begin{equation}
\label{eq:SOodd_box_sum}
\sum_{(i,j)\in\lambda}(b_i+j-\tfrac12)
=\sum_{i=1}^n\sum_{j=1}^{\lambda_i}(b_i+j-\tfrac12)
=\frac12\sum_{i=1}^n\big(\lambda_i^2+2b_i\lambda_i\big)
=\frac12\big(C_2(\lambda)-|\lambda|\big),
\end{equation}
using \Cref{eq:SOodd_casimir_formula}.
Since $\lambda_i^2+\lambda_i\ge 2\lambda_i$ for $\lambda_i\in\N$ and $b_i\ge 0$, we have
$C_2(\lambda)=\sum_i(\lambda_i^2+2b_i\lambda_i+\lambda_i)\ge 2|\lambda|$.
Thus $C_2(\lambda)-|\lambda|\ge \tfrac12 C_2(\lambda)$, so
\begin{equation}
\label{eq:SOodd_box_sum_ge_c2}
\sum_{(i,j)\in\lambda}(b_i+j-\tfrac12)\ \ge\ \frac14\,C_2(\lambda).
\end{equation}
Also $\lambda_1\le \sqrt{C_2(\lambda)}$, hence
$K_{\max}\le (n-1)+\sqrt{C_2(\lambda)}$.
Substitute \Cref{eq:SOodd_box_sum_ge_c2} into \Cref{eq:SOodd_sum_lower} and then into \Cref{eq:SOodd_log_asinh}:
\begin{equation}
\label{eq:SOodd_log_final}
\log a_{\lambda}\ \le\ -\frac14\,\frac{C_2(\lambda)}{t+(n-1)+\sqrt{C_2(\lambda)}}.
\end{equation}
Finally, since $t\le 2\beta/(2n+1)\le \beta$ and $\beta\ge 1$,
$t+(n-1)+\sqrt{C_2}\le \beta+n+\sqrt{C_2}\le (n+1)(\beta+\sqrt{C_2})$.
Therefore \Cref{eq:SOodd_log_final} implies \Cref{eq:SOodd_A1_final} with
$c_{B}(n)=\frac{1}{4(n+1)}$.
\end{proof}

